\documentclass[../../main.tex]{subfiles}
\begin{document}

{\bf [解]}

\begin{figure}[htb]
  \centering
  \begin{tikzpicture}[scale=1.8, >=stealth, every node/.style={font=\small}]

    % --- 1. 座標定義 (例として theta = 120°, |b| = 1.6, |c| = 1.6) ---
    \coordinate (A) at (0, 0);

    % AB の偏角 -115°, AC の偏角 -55° (なす角 theta = 60° ではなく鈍角にするため、
    % AB を -120°、AC を -40° に設定 -> なす角 80°〜120° 程度が見やすい)
    % ここでは手描き図に合わせて AB: -125°, AC: -45° (なす角 80° 程度で描くか、解の 120° に準拠)
    % 解の条件: theta は鈍角 (例えば 110°)
    \def\angB{-135}
    \def\angC{-25}
    \def\lenB{1.8}
    \def\lenC{1.8}

    \coordinate (B) at (\angB:\lenB);
    \coordinate (C) at (\angC:\lenC);

    % 直線 k: A を通り AC に垂直 (偏角 angC + 90° = 65°)
    \def\angK{\angC + 90}
    \coordinate (K_dir) at (\angK:1);

    % 直線 h: A を通り AB に垂直 (偏角 angB + 90° = -45° または -35°)
    \def\angH{\angB + 90}
    \coordinate (H_dir) at (\angH:1);

    % --- 2. 垂足 P, Q および対称点 B' の幾何計算 ---
    % P: B から直線 k (A を通り方向 K_dir) への垂足
    \coordinate (P) at ($ (A)!(B)!($ (A) + (K_dir) $) $);

    % Q: C から直線 h (A を通り方向 H_dir) への垂足
    \coordinate (Q) at ($ (A)!(C)!($ (A) + (H_dir) $) $);

    % B': B の直線 k に関する鏡映点 (B + 2*BP)
    \coordinate (Bprime) at ($ (B)!2.0!(P) $);

    % --- 3. 直線 k と h の描画 ---
    \draw[thick, dashed, gray!80!black]
    ($ (A) - 2.4*(K_dir) $) -- ($ (A) + 2.4*(K_dir) $) node[above right] {$k$};
    \draw[thick, dashed, gray!80!black]
    ($ (A) - 2.2*(H_dir) $) -- ($ (A) + 2.2*(H_dir) $) node[below right] {$h$};

    % --- 4. 三角形の辺 AB, AC ---
    \draw[very thick] (B) -- (A) -- (C);

    % --- 5. 垂線 BP(B') および CQ ---
    % B から P を経て B' への垂線（破線）
    \draw[thick, dashed, blue!80!black] (B) -- (Bprime);
    \draw[thick, dotted, gray!70] (A) -- (Bprime);

    % C から Q への垂線
    \draw[thick, dashed, blue!80!black] (C) -- (Q);
    \draw[thick, dotted, gray!70] (Q) -- ++($ 0.5*(Q) - 0.5*(C) $);

    % --- 6. 直角記号 ---
    % (1) 点 A における k \perp AC
    \coordinate (unitAC) at ($ (A)!0.15!(C) $);
    \coordinate (unitK)  at ($ (A)!0.15!($ (A) + (K_dir) $) $);
    \draw[thick] (unitAC) -- ($ (unitAC) + (unitK) $) -- (unitK);

    % (2) 点 A における h \perp AB
    \coordinate (unitAB) at ($ (A)!0.15!(B) $);
    \coordinate (unitH)  at ($ (A)!0.15!($ (A) - (H_dir) $) $);
    \draw[thick] (unitAB) -- ($ (unitAB) + (unitH) $) -- (unitH);

    % (3) 点 P における垂足直角記号 (BP \perp k)
    \coordinate (uP_B) at ($ (P)!0.15!(B) $);
    \coordinate (uP_K) at ($ (P)!0.15!(A) $);
    \draw[thick, blue!80!black] (uP_B) -- ($ (uP_B) + (uP_K) - (P) $) -- (uP_K);

    % (4) 点 Q における垂足直角記号 (CQ \perp h)
    \coordinate (uQ_C) at ($ (Q)!0.15!(C) $);
    \coordinate (uQ_H) at ($ (Q)!0.15!(A) $);
    \draw[thick, blue!80!black] (uQ_C) -- ($ (uQ_C) + (uQ_H) - (Q) $) -- (uQ_H);

    % --- 7. 頂点・点プロットとラベル ---
    \fill (A) circle (1.2pt) node[above=3pt] {$A$};
    \fill (B) circle (1.4pt) node[below left] {$B$};
    \fill (C) circle (1.4pt) node[below right] {$C$};
    \fill[blue!80!black] (P) circle (1.2pt) node[below left=2pt] {$P$};
    \fill[blue!80!black] (Q) circle (1.2pt) node[below right=2pt] {$Q$};
    \fill[red!80!black] (Bprime) circle (1.3pt) node[above left] {$B'$};

  \end{tikzpicture}
  \caption{直線 $k \perp AC$，$h \perp AB$ および各頂点からの垂線と対称点 $B'$}
  \label{fig:triangle_perpendicular_reflections}
\end{figure}

$B, C$ から $k, h$ に下ろした垂線の足を順に $P, Q$ とおく．
正射影ベクトルの公式から
\begin{align}
  \vec{BP} &= -\dfrac{\vec{b} \cdot \vec{c}}{|\vec{c}|^2} \vec{c} \\
  \vec{CQ} &= -\dfrac{\vec{b} \cdot \vec{c}}{|\vec{b}|^2} \vec{b}
\end{align}
だから，$\vec{BB'} = 2\vec{BP}$，$\vec{CC'} = 2\vec{CQ}$ に注意して
\begin{align}
  \begin{cases}
    \vec{b'} = \vec{b} + 2\vec{BP} = \vec{b} - 2\dfrac{\vec{b} \cdot \vec{c}}{|\vec{c}|^2} \vec{c} \\
    \vec{c'} = \vec{c} + 2\vec{CQ} = -2\dfrac{\vec{b} \cdot \vec{c}}{|\vec{b}|^2} \vec{b} + \vec{c}
  \end{cases} \label{eq:1}
\end{align}
である．

ここで$\angle BAC = \theta \; (0 < \theta < \pi)$ とおく．
題意より$\vec{b}, \vec{c}$ は三角形をなすから一次独立である．
よって\cref{eq:1}および題意の式
\begin{align}
  \begin{cases}
    \vec{b'} = \vec{b} + \vec{c} \\
    \vec{c'} = m\vec{b} + \vec{c}
  \end{cases}
\end{align}
で係数を比較してよく，
\begin{align}
  &
  \begin{cases}
    -2\dfrac{\vec{b} \cdot \vec{c}}{|\vec{c}|^2} = 1 \\
    m &= -2\vec{b} \cdot \vec{c}
  \end{cases}
  \therefore
  &
  \begin{cases}
    -2 \dfrac{|\vec{b}||\vec{c}|\cos\theta}{|\vec{c}|^2} = 1 \\
    m &= -2|\vec{b}||\vec{c}|\cos\theta
  \end{cases}
  \therefore
  &
  \begin{cases}
    -2 \dfrac{\cos\theta}{|\vec{c}|} = 1 \\
    m &= -2|\vec{c}|\cos\theta
  \end{cases} \label{eq:2}
\end{align}
となる．ただし最終行で題意の条件$|\vec{b}|=1$を用いた．

前者から
\begin{align}
  |\vec{c}| = -2\cos\theta
\end{align}
だから，$|\vec{c}| > 0$ より
\begin{align}
  \dfrac{\pi}{2} < \theta < \pi \label{eq:3}
\end{align}
である．このもとで\cref{eq:2}の第二式に代入して
\begin{align}
  m = 4\cos^2\theta \label{eq:4}
\end{align}
となる．$-1 < \cos\theta < 0$ 及び $m \in \mathbb{Z}_{>0}$ より
あり得る$(m,\cos\theta)$の組み合わせは
\begin{align}
  (m, \cos\theta) = \left(1, -\dfrac{1}{2}\right), \left(2, -\dfrac{1}{\sqrt{2}}\right), \left(3, -\dfrac{\sqrt{3}}{2}\right)
\end{align}
である．それぞれのケースで\cref{eq:2}により$|\vec{c}|$を求めて
\begin{align}
  (m, \theta, |\vec{c}|) = \left(1, \dfrac{2}{3}\pi, 1\right), \left(2, \dfrac{3}{4}\pi, \sqrt{2}\right), \left(3, \dfrac{5}{6}\pi, \sqrt{3}\right)
\end{align}
が求める組である．

\end{document}
